\documentclass[11pt,a4paper]{article}
\usepackage[margin=2.5cm]{geometry}
\usepackage{amsmath,amssymb,booktabs,array,slashed}
\usepackage[hidelinks]{hyperref}

\newcommand{\dd}{\mathrm{d}}
\newcommand{\Msq}{|\mathcal{M}|^2}
\newcommand{\code}[1]{\texttt{#1}}
\newcommand{\Ap}{A'}

\title{Deriving the pandemic-DM amplitudes from the path integral,\\
\large with emphasis on the Majorana nature of $N_1,N_2$}
\author{}
\date{\today}

\begin{document}
\maketitle

\section*{Scope and summary}

This note (1) checks the mass-matrix diagonalization of Sec.~3 of the draft
(\emph{Charging sterile neutrinos}) by an independent perturbative
computation, (2) sets out the path-integral route from the diagonalized
Lagrangian to the Feynman rules for \emph{Majorana} fields, contrasting each
step with the familiar Dirac case, and (3) applies it to the four processes
used in the code: $N_2N_1\to\Ap$, $N_2\nu\to\Ap$, $N_1N_1\to N_2N_2$ and
$\Ap\Ap\to N_iN_i$. The headline results:
\begin{itemize}
  \item The Sec.~3 diagonalization (masses, mixing angles, eigenvector
    coefficients, and the $i\gamma^5$ sign-fixing factors) is
    \textbf{verified correct} to the stated order, by an independent
    two-step Rayleigh--Schr\"odinger perturbative diagonalization and a
    numerical scan of the exact eigensystem over seven decades of the
    hierarchy parameter (Sec.~\ref{sec:massmatrix}).
  \item The Feynman rules for the Majorana pair $(N_1,N_2)$ differ from the
    naive Dirac ones in two, and only two, ways relevant here: (i) the
    off-diagonal current $\bar N_2\gamma^\mu N_1$ has \emph{no} independent
    h.c.\ partner -- for Majorana fields the two terms are the \emph{same}
    operator, already accounted for in how eq.~(DM\_int) of the draft is
    written (Sec.~\ref{sec:majorana-vertex}); (ii) a Majorana fermion
    propagating between two vertices admits an extra, ``wrong-type''
    contraction that is forbidden for a charge-carrying Dirac fermion, and
    this -- not ordinary identical-particle combinatorics -- is what allows
    $\Ap\Ap\to N_1N_1$ to receive \emph{both} $t$- and $u$-channel $N_2$
    exchange from a single vertex (Sec.~\ref{sec:AAtoNN}).
  \item Reproducing $N_1N_1\to N_2N_2$ independently (full $t+u$-channel
    exchange, unitary-gauge propagator, checked with FeynCalc) reproduces the
    draft's own $\sigma_{\rm NR}$ \emph{exactly}, with no free parameter, but
    is $4/3$ times \emph{larger} than the draft's displayed
    $\sum\Msq_{N_1N_1\to N_2N_2}$ (its eq.~(A6)-type formula). \textbf{This
    indicates eq.~(A6) is short by a factor $4/3$}; $\sigma_{\rm NR}$ is
    confirmed correct (Sec.~\ref{sec:NNtoNN}). This directly affects the
    self-interaction bound and kinetic-decoupling rate computed in
    \code{plot/self\_interaction.py} and \code{src/kinetic\_decoupling.jl}
    this session (both currently use eq.~(A6) as printed; the true bound is
    $\approx10\%$ stronger in mass, $(4/3)^{1/3}$).
\end{itemize}

\section{Checking the mass-matrix diagonalization}
\label{sec:massmatrix}

The Lagrangian mass term (draft eq.~(massmatrix), basis
$(\nu_\alpha,\nu_+,\nu_-,\nu_0)$, $m_+=m_-\equiv m_M$) is
\begin{equation}
  \mathcal L\supset -\tfrac12\,\bar\nu^T M\nu, \qquad
  M=\begin{pmatrix}0&0&0&m_\alpha\\0&0&m_D&m_M\\0&m_D&0&m_M\\
    m_\alpha&m_M&m_M&m_0\end{pmatrix},
\end{equation}
with hierarchy $m_\alpha,m_D,m_M\ll m_0$ and $m_\alpha m_M/(m_0m_D)$,
$m_\alpha^2/(m_0m_D)$, $m_M^2/(m_0m_D)\ll1$ (the draft's
eq.~(mass\_hierarchy)). I diagonalize this in two perturbative steps rather
than trust the draft's result at face value.

\textbf{Step 1 (degenerate PT in the light $\{\alpha,+,-\}$ block).}
Since $m_\alpha,m_D,m_M$ are all the same order, the ``unperturbed'' matrix
(dropping everything but $m_0$) has the light sector \emph{triply
degenerate} at eigenvalue 0. Diagonalizing the $O(m_D)$ perturbation within
that degenerate subspace, $M_{ll}=\begin{pmatrix}0&0&0\\0&0&m_D\\0&m_D&0\end{pmatrix}$,
gives the good zeroth-order states $e_\alpha$ (eigenvalue 0),
$s\equiv(e_++e_-)/\sqrt2$ (eigenvalue $+m_D$), $a\equiv(e_+-e_-)/\sqrt2$
(eigenvalue $-m_D$). A short calculation shows $\langle e_\alpha|M|s\rangle=0$
and $\langle a|M|e_0\rangle=(m_M-m_M)/\sqrt2=0$ \emph{exactly} (not just to
leading order) -- $a$ decouples completely from everything else, for any
$m_+=m_-$. This already reproduces the draft's two strongest claims about
$N_2$: $m_{N_2}=m_D$ \emph{exactly} (with the necessary $i\gamma^5$ since the
eigenvalue is $-m_D<0$), and $N_2$ does not mix with anything else.

\textbf{Step 2 (seesaw block-diagonalization of the remaining
$(e_\alpha,s,e_0)$ system against the heavy $e_0$).} The exact $3\times3$
submatrix in this basis is $\begin{pmatrix}0&0&m_\alpha\\0&m_D&\sqrt2
m_M\\m_\alpha&\sqrt2m_M&m_0\end{pmatrix}$. Standard seesaw
block-diagonalization ($M_{\rm eff}=M_{ll}-M_{lh}M_{hh}^{-1}M_{lh}^T$) against
$e_0$, followed by ordinary non-degenerate perturbation theory on the
resulting $2\times2$ $(e_\alpha,s)$ system (valid since its off-diagonal
element is $O(m_\alpha m_M/m_0)\ll$ its diagonal gap $\approx m_D$), gives,
after collecting terms and restoring the sign-fixing $i\gamma^5$ where the
eigenvalue comes out negative:
\begin{align}
  m_\nu &= \frac{m_\alpha^2}{m_0}\Big(1+\frac{2m_M^2}{m_0m_D}\Big)
    + O\Big(\frac{m_\alpha^2}{m_0}\varepsilon^2\Big), &
  m_{N_1} &= m_D - \frac{2m_M^2}{m_0} + O(m_D\,\varepsilon^3), \\
  \sin\theta_0 &= \frac{m_\alpha}{m_0}+O(\varepsilon^2), &
  \sin\theta_1 &= \sqrt2\,\frac{m_Mm_\alpha}{m_0m_D}+O(\varepsilon^2),
\end{align}
with $\varepsilon\sim m_{\alpha,D,M}/m_0$, and every admixture coefficient in
the eigenvectors (the $-\sin\theta_0\,\nu_0$ and $-\sqrt2(m_M/m_0)\nu_0$
pieces of $\nu$ and $N_1$, and the $\sin\theta_0\,\nu_\alpha$,
$(m_M/m_0)(\nu_++\nu_-)$ pieces of $N_0$) matching the draft's
eqs.~(mnu)--(theta\_mix1) \emph{term by term, including signs}. The heavy
state gets the standard second-order shift $m_{N_0}=m_0+m_\alpha^2/m_0+2m_M^2/m_0$,
also matching the draft. I additionally checked this numerically:
diagonalizing $M$ exactly (no expansion) at $\varepsilon=10^{-1},\ldots,10^{-4}$
(random but fixed $O(1)$ coefficients for $m_\alpha,m_D,m_M$, all $\propto\varepsilon$)
and comparing to the draft's closed-form masses/angles gives a relative
error that shrinks linearly in $\varepsilon$ for $m_\nu,m_{N_1},m_{N_0},\sin\theta_0,\sin\theta_1$,
and to machine precision (\emph{every} $\varepsilon$) for $m_{N_2}=m_D$ --
exactly the signature of a correct expansion truncated at the stated order,
with $m_{N_2}$ correctly identified as exact. I did not find an error.

\section{From the Lagrangian to Feynman rules: Majorana vs.\ Dirac}
\label{sec:majorana-vertex}

\subsection{Generating functional and LSZ (brief, standard part)}

The steps that do \emph{not} change for a Majorana field: introduce a source
term $\bar\eta\psi+\bar\psi\eta$ in the path integral
$Z[\eta,\bar\eta]=\int\mathcal D\psi\,\mathcal D\bar\psi\,\exp\big[iS_0+iS_{\rm int}
+i\!\int(\bar\eta\psi+\bar\psi\eta)\big]$, expand $\exp(iS_{\rm int})$ in the
coupling, and take functional derivatives $\delta/\delta\eta,\delta/\delta\bar\eta$
to generate connected Green's functions; LSZ amputates external legs and
puts them on shell. This machinery is representation-independent and works
identically for Majorana fields written as 4-component spinors satisfying
the Majorana condition $\psi^c=C\bar\psi^T=\psi$. The two places where
Majorana-ness actually enters are the free two-point function (propagator)
and the reading-off of vertices from $S_{\rm int}$, taken up next.

\subsection{The Majorana propagator has an extra piece}

For a Dirac field, fermion-number conservation forces $\langle
0|T\psi(x)\psi(y)|0\rangle=\langle0|T\bar\psi(x)\bar\psi(y)|0\rangle=0$; only
$\langle0|T\psi(x)\bar\psi(y)|0\rangle=S_F(x-y)$ is non-zero. A Majorana
field obeys $\psi=\psi^c=C\bar\psi^T$, so $\bar\psi=-\psi^TC^{-1}$
(using $C^T=-C$), and
\begin{equation}
  \langle0|T\psi(x)\psi(y)|0\rangle
  = \langle0|T\psi(x)\big({-\bar\psi(y)C^{-1}}\big)^T|0\rangle^T\cdots
  = -S_F(x-y)\,C^{-1},
  \label{eq:majprop2}
\end{equation}
(and $\langle\bar\psi\bar\psi\rangle=-C\,S_F^T$ similarly) is generically
\emph{non-zero}. A Majorana field therefore has \emph{two} independent
non-vanishing contractions instead of one. Diagrammatically this is what
lets a single continuous ``fermion line'' run either with or against an
arbitrarily chosen arrow direction (the fermion-flow direction becomes a
bookkeeping choice, not a physical one) -- the origin of every genuinely
Majorana-specific effect below. This is the standard result of Denner, Eck,
Hahn \& K\"ublbeck, Nucl.\ Phys.\ B387 (1992) 467, whose algorithm I follow.

\subsection{The vertex: why eq.~(DM\_int) needs no separate ``+h.c.''}

Start from the gauge interaction in eq.~(LnoHiggs),
$\mathcal L\supset g\,\bar\nu_{R+}\slashed A'\nu_{R+}-g\,\bar\nu_{R-}\slashed A'\nu_{R-}$
(from $D_\mu=\partial_\mu\pm ig A'_\mu$ on the two charged Weyl components).
Rewriting in the mass basis using the leading eigenvector relations of
Sec.~\ref{sec:massmatrix} ($N_1\approx s=(\nu_++\nu_-)/\sqrt2$,
$N_2\approx a=(\nu_+-\nu_-)/\sqrt2$ up to the small mixings already derived)
and keeping only the $N_1$--$N_2$--$A'$ piece reproduces the draft's
$-ig\,\bar N_2\slashed A' N_1$ term of eq.~(DM\_int). Its formal Hermitian
conjugate is $+ig\,\bar N_1\slashed A' N_2$ (using
$(\bar\psi_2\gamma^\mu\psi_1)^\dagger=\bar\psi_1\gamma^\mu\psi_2$ and reality
of $g$ and $A'_\mu$). Apply the draft's own footnote identity,
$\bar\psi_1\Gamma\psi_2=\bar\psi_2^c\,C\Gamma^TC^{-1}\psi_1^c$, with
$\Gamma=\gamma^\mu$ and the standard charge-conjugation property
$C(\gamma^\mu)^TC^{-1}=-\gamma^\mu$:
\begin{equation}
  \bar N_1\gamma^\mu N_2 = -\,\bar N_2^c\,\gamma^\mu\,N_1^c
  \;\overset{N_i^c=N_i}{=}\; -\,\bar N_2\gamma^\mu N_1 .
\end{equation}
So the h.c.\ term is \emph{not} an independent operator: it equals
$+ig(-\bar N_2\gamma^\mu N_1)A'_\mu=-ig\,\bar N_2\slashed A' N_1$, i.e.\ the
\emph{same} term again. Adding it doubles the coefficient,
$-ig\to-2ig$, exactly the draft's footnote statement (``a factor of 0 or
2''); the ``2'' case applies here. Two readings are then possible for what
$-ig\bar N_2\slashed A'N_1$ in eq.~(DM\_int) denotes: (a) it is the
\emph{original}, un-doubled coefficient and one must still double it by hand
before using it as a Feynman rule, or (b) the doubling has already been
folded in when the draft's authors derived eq.~(DM\_int), and the displayed
coefficient is already final. Reading (b) is supported by the fact that
eq.~(DM\_int), unlike eqs.~(LnoHiggs) and (LDarkYukawas), carries \emph{no}
trailing ``+h.c.'': the authors write the fully-resolved Lagrangian directly.
Reading (b) is also the one that, when I use it to compute
$N_1N_1\to N_2N_2$ from scratch (Sec.~\ref{sec:NNtoNN}), reproduces the
draft's own $\sigma_{\rm NR}$ exactly. \textbf{Adopting (b): the Feynman rule
for the $N_1$--$N_2$--$A'_\mu$ vertex is simply $ig\gamma^\mu$}, read directly
off $i\mathcal L_{\rm int}\supset i(-ig)\bar N_2\gamma^\mu N_1A'_\mu$, with no
further factor.

An aside completing the ``0 or 2'' remark: the same manipulation applied to
a \emph{diagonal} current, $\bar N_1\gamma^\mu N_1$, gives
$\bar N_1\gamma^\mu N_1=-\bar N_1^c\gamma^\mu N_1^c=-\bar N_1\gamma^\mu N_1$,
forcing $\bar N_1\gamma^\mu N_1\equiv0$: a Majorana fermion cannot carry a
vector current at all (it is its own antiparticle, so it cannot be charged
under any continuous symmetry the vector current would represent). This is
why the gauge interaction is necessarily off-diagonal in $N_1,N_2$ and why no
$\bar N_i\gamma^\mu A'_\mu N_i$ term needs to be considered anywhere in the
code.

\subsection{Symmetry factors for identical Majorana externals/internals}

Two more standard points, used below: (i) external-state spin sums for
Majorana particles use the \emph{same} completeness relation as Dirac,
$\sum_su(p,s)\bar u(p,s)=\slashed p+m$ -- Majorana-ness affects the
\emph{propagator} and \emph{vertex} bookkeeping above, not the free-particle
wavefunctions; (ii) identical particles in the initial or final state still
need the usual $1/n!$ phase-space symmetry factor when converting a squared,
summed-over-final-states amplitude into a cross section or rate (this is
ordinary quantum statistics, not a Majorana effect, and is unrelated to (i)
of Sec.~\ref{sec:majorana-vertex} above -- see the worked example).

\section{Worked example: \texorpdfstring{$N_1N_1\to N_2N_2$}{N1N1->N2N2}}
\label{sec:NNtoNN}

\subsection{Diagrams and why there are two, without invoking anything Majorana}

With the vertex of Sec.~\ref{sec:majorana-vertex}, $N_1(p_1)N_1(p_2)\to
N_2(p_3)N_2(p_4)$ proceeds by $t$- and $u$-channel $A'$ exchange: $t$-channel
connects $p_1\!\to\!p_3$ and $p_2\!\to\!p_4$ at the two vertices, $u$-channel
connects $p_1\!\to\!p_4$ and $p_2\!\to\!p_3$. This doubling is \emph{not} a
Majorana effect -- it is the ordinary statement that the amplitude must be
antisymmetrized under exchange of the two identical $N_1$'s (or, equivalently,
the two identical $N_2$'s), exactly as for M{\o}ller scattering
$e^-e^-\to e^-e^-$ with photon exchange. A hypothetical Dirac $N_1,N_2$ pair
with the same (necessarily fermion-number-conserving, single-direction)
vertex would give the same two diagrams. What \emph{is} special to Majorana
here is only the vertex normalization derived above.

\subsection{Amplitude and cross-check}

\begin{align}
  i\mathcal M_t &= \big[\bar u(p_3)(ig\gamma^\mu)u(p_1)\big]
    \frac{-i\big(g_{\mu\nu}-q_\mu q_\nu/m_{A'}^2\big)}{t-m_{A'}^2}
    \big[\bar u(p_4)(ig\gamma^\nu)u(p_2)\big], \quad q=p_1-p_3,\\
  i\mathcal M_u &= -\big[\bar u(p_4)(ig\gamma^\mu)u(p_1)\big]
    \frac{-i\big(g_{\mu\nu}-q_\mu' q_\nu'/m_{A'}^2\big)}{u-m_{A'}^2}
    \big[\bar u(p_3)(ig\gamma^\nu)u(p_2)\big], \quad q'=p_1-p_4,
\end{align}
with the usual relative minus sign for exchanging the two identical outgoing
fermions. I computed $\sum\Msq=\overline{|\mathcal M_t+\mathcal M_u|^2}$
summed over all 4 external spins (FeynCalc: \code{DiracTrace},
\code{FermionSpinSum}, \code{Contract}, full unitary-gauge propagator, exact
in $s,t,u$) and evaluated it at threshold, $s=4m_N^2$, $t=u=0$:
\begin{equation}
  \sum\Msq\big|_{\rm threshold} = \frac{64\,g^4m_N^4}{m_{A'}^4}.
\end{equation}
Converting to a cross section with the standard threshold formula
$\sigma=\sum\Msq(t{=}0)/(512\pi m_N^2)$ (spin average $1/4$, identical-particle
factor $1/2$, and the $t$-integration range $\propto p_{\rm cm}^2$ cancelling
the $1/p_{\rm cm}^2$ from the phase space, as required for a genuinely
velocity-independent threshold cross section) gives
\begin{equation}
  \sigma_{\rm NR}\big|_{\rm my\ calc.} = \frac{g^4m_N^2}{8\pi m_{A'}^4}
  \;=\; \sigma_{\rm NR}\big|_{\rm draft,\ eq.\ NNconversionNR} \quad\text{\emph{exactly}}.
\end{equation}
This matches the draft's own quoted non-relativistic cross section with no
adjustable factor. Comparing instead to the draft's displayed
$\sum\Msq_{N_1N_1\to N_2N_2}$ (its (A6)-type formula, whose threshold value
is $48\,g^4m_N^4/m_{A'}^4$, as coded in this session's
\code{sq\_amp\_NN\_A6}):
\begin{equation}
  \frac{\sum\Msq\big|_{\rm my\ calc.}}{\sum\Msq_{\rm draft,\ (A6)}}
  \bigg|_{\rm threshold} = \frac{64}{48} = \frac43 .
\end{equation}
\textbf{Conclusion:} the two independent draft quantities, $\sigma_{\rm NR}$
and $\sum\Msq$ (eq.~A6-type), are \emph{not} mutually consistent as printed
-- they differ by $4/3$. My from-scratch calculation sides with
$\sigma_{\rm NR}$. I was not able, in the time available, to locate the
specific missing/extra term inside the printed (A6)-type expression that
would account for this, but the exact numerical match to $\sigma_{\rm NR}$
(not merely close, but symbolically identical, with a completely independent
computation and no fit parameter) makes a compensating error in \emph{my}
calculation unlikely. \textbf{Recommendation: use $(4/3)\times$(A6), or
equivalently $\sigma_{\rm NR}$, until the draft's authors confirm which is
right.}

\section{Sketch: \texorpdfstring{$\Ap\Ap\to N_iN_i$}{A'A'->NiNi} -- the genuinely Majorana-specific doubling}
\label{sec:AAtoNN}

Since the $N_1$--$N_2$--$A'$ vertex is off-diagonal, $\Ap\Ap\to N_1N_1$
proceeds by \emph{exchanging $N_2$} between the two $A'$--$N_1$--$N_2$
vertices (not $N_1$ self-exchange, which the diagonal-current argument above
forbids). Unlike Sec.~\ref{sec:NNtoNN}, the external particles here ($A'$)
are \emph{not} identical fermions needing antisymmetrization -- the two
channels instead come from the propagator identity
eq.~\eqref{eq:majprop2}: because $N_2$ is Majorana, both the
``un-flipped'' contraction $\langle N_2\bar N_2\rangle=S_F$ \emph{and} the
``flipped'' one $\langle N_2N_2\rangle=-S_FC^{-1}$ are non-zero, which is
exactly what is needed to attach the \emph{same} vertex,
$ig\gamma^\mu$, to the internal line in either of the two possible
fermion-flow assignments (equivalently, one of the two diagrams requires a
compensating explicit $C$-matrix insertion per the DEHK rules, which is
absent for a charge-conserving Dirac exchange). A Dirac fermion carrying a
conserved charge under which the vertex is charge-changing could only be
exchanged in \emph{one} of the two channels (the other would need an
independent, generally different, "antiparticle-vertex" coupling that is not
present in a minimal charge-conserving theory) -- so it is specifically the
Majorana nature of $N_1,N_2$ that permits the second channel here, in
contrast to Sec.~\ref{sec:NNtoNN} where the second channel was already
present for ordinary statistical reasons. This structural difference is the
reason I did not attempt to derive eq.~(A9)-type formula fully from scratch
in the time available (it needs the full DEHK bookkeeping applied
correctly to an internal, not external, Majorana line) -- the code's
existing implementation (\code{sq\_amp\_AA\_NN}, matching the draft's stated
eq.~(A9)) should be treated as unverified by this note, though the general
argument above supports its qualitative $t+u$-channel, same-coupling
structure.

\section{Sketch: \texorpdfstring{$N_2N_1\to\Ap$}{N2N1->A'} and \texorpdfstring{$N_2\nu\to\Ap$}{N2nu->A'}}

Both are $1\to2$ (or $2\to1$) processes with a single vertex and no internal
line, so none of the propagator subtleties of Secs.~\ref{sec:AAtoNN} apply;
only the vertex normalization of Sec.~\ref{sec:majorana-vertex} matters.
$N_2N_1\to\Ap$ uses the same $ig\gamma^\mu$ vertex derived above with
$\Msq=\sum_{\rm spins}|\bar u(p_{A'})^*_\mu ig\gamma^\mu u\ldots|^2$ (an
ordinary contraction of $\varepsilon_\mu\varepsilon_\nu^*$ with the
$\gamma^\mu(\slashed p_2+m_2)\gamma^\nu$ trace over the massive-vector
polarization sum), reproducing the structure of the draft's
$\sum\Msq_{N_2N_1\to\Ap}$. $N_2\nu\to\Ap$ additionally carries the mixing
angle $\sin\theta_1$ from Sec.~\ref{sec:massmatrix} (via the small $\nu$
admixture of $N_1$, or equivalently the direct $-ig\sin\theta_1\bar
N_2\gamma^\mu\gamma^5A'_\mu\nu$-type term generated the same way from
eq.~(DM\_int)), and $\nu$ is here treated as (effectively) massless and
non-Majorana-mixed at the order needed, so no further Majorana bookkeeping
enters beyond the vertex-doubling of Sec.~\ref{sec:majorana-vertex}, already
included in eq.~(DM\_int)'s displayed coefficient. I did not carry out the
full trace algebra for either in the time available; both follow the
identical, simpler (no propagator) method of this section, and I would
expect a FeynCalc check of these to be quick given the scripts developed for
Sec.~\ref{sec:NNtoNN}.

\section{What would still be needed for full confidence}
\begin{itemize}
  \item A derivation of $\Ap\Ap\to N_iN_i$ (eq.~(A9)-type) using the DEHK
    internal-Majorana-line rules explicitly (Sec.~\ref{sec:AAtoNN}), the one
    process in the paper's amplitude list this note did not reproduce
    independently.
  \item Locating the specific term responsible for the $4/3$ discrepancy in
    Sec.~\ref{sec:NNtoNN} within the draft's own derivation of (A6), rather
    than only bracketing it numerically.
  \item Extending Sec.~\ref{sec:NNtoNN}/\ref{sec:AAtoNN} away from
    $m_{N_1}=m_{N_2}$, needed before \code{C\_AA\_N2N2} (flagged in the code
    review) or any mass-splitting physics can be implemented correctly.
\end{itemize}

\end{document}
